\section{Cyclotomic derivatives at the initial and auxiliary primes}
\label{sec:minus-derivatives}

We establish the two local relations for \(\nu_L\) used in
Section~\ref{sec:minus-characterwise}. Fix \(M=\ell^m\), \(m\ge1\),
\(\ell\) odd, \(P=N_0M\) with \(\ell\nmid N_0\), and a prime
\(r\equiv1\bmod P\). Put \(K_n=\mathbb Q(\mu_{Pn})\) and
\(H=G_{K_1}\). Each index \(n\) is the product of a list \(L\) of
distinct admissible primes, in the sense of the preceding section.
Let \(\alpha\) be the character there and \(\beta=\kappa\alpha^{-1}\).
On \(H\), \(\alpha\equiv1\bmod M\). Choose
\(\delta\) with \(\alpha(\delta)-\kappa(\delta)\) a unit, and an
integer \(a\) giving the action of \(\delta\) on \(\mu_M\).
Thus \(\alpha(\delta)-a\) is a unit. Write \(T(v)=\delta(v)/v^a\).
The map \(T\) is multiplicative and kills \(\mu_M\).

\subsection{The input cyclotomic elements and the derivatives}

The cyclotomic input supplies \(X_n\in K_n^\times\), units away
from \(r\), and integers \(s(n,q)\ge1\) prime to
\(\ell\) for each admissible extension \(n\mapsto qn\).
For \(\tau_q\) acting as a generator on \(\mu_q\) and fixing the
prime-to-\(q\) cyclotomic factors, and \(\phi\) acting on \(K_n\)
by the exponent \(q\), the input says
\[
  \operatorname N_q X_{qn}
    =\eta\,(X_n/\phi^{-1}X_n)^{s(n,q)},\qquad
  X_{qn}\equiv\eta'\,(\phi^{-1}X_n)^{s(n,q)}\quad\bmod\mathfrak q
\]
at all places \(\mathfrak q\mid q\). Here \(\operatorname N_q\)
denotes the norm \(K_{qn}/K_n\). Both error factors are roots
of unity of order dividing \(2Pqn\), and the same \(\eta'\) works
for all these places. At a fixed place \(A_0\mid r\), normalized by
\(v_{A_0}(r)=1\), the prescribed initial divisor is
\(v_{A_0}(\sigma_t X_1)=h_b(t)\) for every lift \(\sigma_t\) of
\(t\in(\mathbb Z/P\mathbb Z)^\times\), with \(h_b\) as in the
preceding section. This is the input called \texttt{laxEx}. Its
construction uses cyclotomic ideal factorizations, relative norms
and congruences of chosen generators.

Put \(x_n=T(X_n)^{2N_0n}\) and \(c(n,q)=q\,s(n,q)\).
If \(\eta^{2Pqn}=1\), then \(\eta^{2N_0qn}\in\mu_M\), so
\(T(\eta)^{2N_0qn}=1\). Norm and the automorphism \(\delta\) commute
on these abelian cyclotomic fields. Conjugating a place above \(q\)
gives another place above \(q\), so \(T\) preserves the congruences
with the indicated transformed error. Applying \(T\) and raising
to \(2N_0qn\) gives
\begin{equation}
 \operatorname N_q x_{qn}=(x_n/\phi^{-1}x_n)^{c(n,q)},\qquad
 x_{qn}\equiv(\phi^{-1}x_n)^{c(n,q)}\pmod{\mathfrak q}.
 \label{eq:deriv-norm}
\end{equation}
The exponents here explain the factor \(q\) in the final transition
unit. Since \(q\equiv1\bmod P\), it is prime to \(\ell\).

For \(v\in K_{qn}^\times\) set
\(D_q(v)=\prod_{i=0}^{q-2}(\tau_q^iv)^i\). Let \(D_L\) be the
composition over \(L\) and set \(y_L=D_L(x_n)\). All these
operators commute with automorphisms on the cyclotomic field
and with each other. Telescoping the consecutive powers and using
\(\tau_q^{q-1}=1\) there gives
\begin{equation}
 \frac{\tau_q(D_q v)}{D_q v}
       =\frac{v^{q-1}}{\operatorname N_q v}.
 \label{eq:deriv-telescoping}
\end{equation}
For every \(g\in H\), \(g(y_L)/y_L\) is an \(M\)-th power in
\(K_n^\times\). To verify this, use induction on \(|L|\).
The empty-list element belongs to \(K_1\). For \(q\in L\) write
\(n=qn'\) and \(L'=L\setminus\{q\}\). Equations \eqref{eq:deriv-norm} and \eqref{eq:deriv-telescoping} give
\[
 \frac{\tau_q(y_L)}{y_L}
  =D_{L'}(x_n)^{q-1}
       \left(\frac{\phi^{-1}y_{L'}}{y_{L'}}\right)^{c(n',q)} .
\]
The first factor is an \(M\)-th power because \(M\mid q-1\).
For the second use the induction hypothesis since \(\phi^{-1}\)
fixes \(K_1\). The restrictions of all these \(\tau_q\) generate
\(\operatorname{Gal}(K_n/K_1)\) by the Chinese remainder theorem.
The elements \(g\) satisfying the assertion form a subgroup, since
\(K_n\) is Galois and the \(M\)-th powers are preserved by every
automorphism. This proves the assertion.

\subsection{The Kummer classes and their projection}

Choose a primitive \(M\)-th root \(\zeta\), and \(Y\) with \(Y^M=y_L\).
For \(g\in H\) take \(v_g\in K_n\) with \(v_g^M=g(y_L)/y_L\).
Define \(\kappa_L(g)\in\mathbb Z/M\) by
\begin{equation}
 \zeta^{\kappa_L(g)}
       =\frac{g(TY)}{TY\,T(v_g)} .
 \label{eq:deriv-kummer}
\end{equation}
The right side has \(M\)-th power one since \(g\) and \(\delta\)
commute on \(K_n\). Changing either root changes it by an element of
\(\mu_M\), which \(T\) kills. Taking
\(v_{gh}=g(v_h)v_g\) proves additivity, since \(T(gv_h)=g(Tv_h)\)
and \(g\) fixes \(\mu_M\). Local constancy follows by restricting to
a finite extension containing \(Y,\delta Y\) and the finitely many
roots \(v_g\) required for the finite possible actions on \(K_n\).
At a prime \(p\nmid Prn\), inertia fixes \(K_n\). The element \(y_L\)
is a \(p\)-adic unit and \(p\nmid M\), so it has an \(M\)-th root
in the local unramified extension, by reduction and Hensel's lemma.
The same holds for \(\delta y_L\) at this place. Thus the chosen
roots are unramified there, and \eqref{eq:deriv-kummer} gives zero on inertia.

Let \(w\) be a generator of the cyclic module \(W[M]\). Using the
lifts fixed in the preceding chapter, set
\begin{equation}
    \nu_L(h)=\sum_t \kappa_L(\sigma_t^{-1}h\sigma_t)
                                 \,\beta(\sigma_t)w,\qquad h\in H .
 \label{eq:deriv-projection}
\end{equation}
This is independent of representatives for the residues of
\(\kappa_L\), since \(Mw=0\). It is locally constant, additive on
\(H\), killed by \(M\), and has the inertia vanishing just established.
Conjugation by \(g\in G_{\mathbb Q}\) multiplies \eqref{eq:deriv-projection} by \(\beta(g)\):
reindex the sum through the action of \(g\) modulo \(H\).
Conjugation by elements of \(H\) preserves the additive homomorphism
\(\kappa_L|_H\), and \(\beta|_H\equiv1\bmod M\).

\subsection{The initial value at \(r\)}

Take \(\rho^M=r\) and choose \(t_0\in I_{A_0}\) with
\(t_0(\rho)=\zeta\rho\). Such an element exists since \(r\nmid M\)
and the \(M\)-th root of the uniformizer generates a totally tamely
ramified extension over the unramified field containing \(\mu_M\).
In the empty-list case \(y=x_1\) belongs to \(K_1\). Write
\(K(t)=v_{A_0}(\sigma_t T(y))\); these are integers. Let \(Y^M=y\)
and \(W_t=\sigma_t T(Y)\). Then \(v_{A_0}(W_t)=K(t)/M\), so
\(W_t/\rho^{K(t)}\) is a unit. The \(M\)-th power of this unit is
fixed by \(t_0\), since \(T(y)\) and its conjugates belong to the
unramified cyclotomic field at \(r\). Its quotient under \(t_0\) is
therefore an \(M\)-th root of unity reducing to one. It equals one,
as \(r\nmid M\). Thus \(t_0(W_t)/W_t=\zeta^{K(t)}\). In \eqref{eq:deriv-kummer} the
element \(v_g\) for \(g=\sigma_t^{-1}t_0\sigma_t\) has \(M\)-th
power one. Applying \(\sigma_t\) to \eqref{eq:deriv-kummer} gives
\[
       [t]\kappa_{\varnothing}(\sigma_t^{-1}t_0\sigma_t)
                         =K(t)\pmod M .
\]

Define \(\mathcal Z(k)=\sum_t k(t)\alpha(\sigma_t)^{-1}\).
Because \(\kappa(\sigma_t)\equiv[t]\bmod M\) for the cyclotomic
character \(\kappa\), equation \eqref{eq:deriv-projection} gives
\(\nu_{\varnothing}(t_0)=\mathcal Z(K)w\).
If \(\delta\) corresponds to \(d\in(\mathbb Z/P\mathbb Z)^\times\),
the operator \(T\) transforms the divisor vector \(k(t)\) into
\(k(td)-a k(t)\). Permuting \(t\mapsto td\) shows that \(\mathcal Z\)
of this new vector equals
\((\alpha(\delta)-a)\mathcal Z(k)\) modulo \(M\), since
\(\alpha|_H\equiv1\bmod M\). Now \(T(y)=T(T(X_1)^{2N_0})\).
Apply this transformation twice and the multiplicity \(2N_0\) once,
starting from \(h_b(t)\). As \(z=\mathcal Z(h_b)\),
\begin{equation}
      \nu_{\varnothing}(t_0)
             =2N_0(\alpha(\delta)-a)^2 zw .
 \label{eq:deriv-initial}
\end{equation}
The multiplier is a unit because \(\ell\nmid2N_0\) and
\(\alpha(\delta)-a\) is a unit. This proves the initial relation with
its actual coefficient.

\subsection{The transition from Frobenius to inertia}

Let \(q\notin L\) be admissible, \(n=\prod_{p\in L}p\),
\(c=c(n,q)\) and \(k=(q-1)/M\). At a fixed place over \(q\) take
an arithmetic Frobenius \(\phi\) and \(t_q\) in inertia whose action
on \(\mu_q\) agrees with \(\tau_q\). Cyclotomic total ramification
at \(q\) supplies \(t_q\). The actions of \(t_q\) and \(\tau_q\)
agree on \(K_{qn}\) since both fix \(K_n\). Put
\[
       y=D_L(x_n),\qquad u=D_L(x_{qn}),\qquad y'=D_q u .
\]
At \(q\) all these are units. Applying \(D_L\) to \eqref{eq:deriv-norm} gives
\(\operatorname N_q u=(y/\phi^{-1}y)^c\) and
\(u\equiv(\phi^{-1}y)^c\) at every place over \(q\).
Take \(Y^M=y'\), \(V^M=y\), and
\(b_0\in K_n\) with \(b_0^M=(\phi^{-1}y)/y\), using the
\(M\)-th power assertion above and \(\phi\in H\).
By \eqref{eq:deriv-telescoping},
\[
       (u^k b_0^c)^M=t_q(y')/y'.
\]
Thus \eqref{eq:deriv-kummer} computes \(\kappa_{qL}(t_q)\) with \(u^k b_0^c\) as its
adjusting root. At \(\phi^{-1}\), the corresponding root for
\(\kappa_L\) is \(b_0\). Dividing the two formulas gives the identity
\begin{equation}
 \frac{\zeta^{\kappa_{qL}(t_q)}}{(\zeta^{\kappa_L(\phi^{-1})})^c}
  =\frac{t_q(TY)}{TY}
     \left(\frac{TV}{(\phi^{-1}TV)^q}\right)^c
      \bigg/ T\!\left(\frac{u}{(\phi^{-1}y)^c}\right)^k .
 \label{eq:deriv-transition-id}
\end{equation}
For the algebra here, \((\phi^{-1}TV)^M=T(\phi^{-1}y)\) since
\(\delta\) and \(\phi\) commute on \(K_n\), and \(q=1+kM\).
The factors \(T(b_0)^c\) in the two expressions cancel.

Every factor on the right of \eqref{eq:deriv-transition-id} reduces to one. For the first use
inertia on the unit \(TY\). For the second use that Frobenius acts
as the \(q\)-power on reductions, applied to the unit
\(\phi^{-1}TV\). For the denominator use the congruence for \(u\)
at every place over \(q\), which also applies at the place conjugate
by \(\delta\). The left side is an \(M\)-th root of unity and \(q\nmid
M\), so injectivity of reduction on \(\mu_M\) makes it one.
Additivity on \(H\) gives
\[
            \kappa_{qL}(t_q)=-c\,\kappa_L(\phi)\pmod M .
\]
This computation applies to
\((\sigma_t^{-1}t_q\sigma_t,\sigma_t^{-1}\phi\sigma_t)\) at the
conjugate place: the new inertia element still agrees with
\(\tau_q\) on \(\mu_q\), because the cyclotomic Galois group is
abelian. Sum the congruences with the weights in \eqref{eq:deriv-projection}. The result is
\begin{equation}
                 \nu_{qL}(t_q)=-q\,s(n,q)\,\nu_L(\phi).
 \label{eq:deriv-transition}
\end{equation}
The factor is a unit because \(\ell\nmid q s(n,q)\). This proves
the transition used in the length induction.

Section~\ref{sec:minus-cyclotomic-elements} constructs the input
\(X_n\), with the prescribed divisor and the same root of unity in
the congruence at every prime above \(q\).
\footnote{\href{https://github.com/ChristianKudeba/alpoge-rh-verification/blob/alpoge-source-7cbf7a0c/src/Development/AutomorphicArithmetic/AutomorphicVarieties/PrimitiveFrobeniusFixed/PlaceEmb.lean}{Cyclotomic elements and divisor data},
\href{https://github.com/ChristianKudeba/alpoge-rh-verification/blob/alpoge-source-7cbf7a0c/src/Development/AutomorphicArithmetic/AutomorphicVarieties/PrimitiveFrobeniusFixed/SysOfLax.lean}{removal of root-of-unity factors},
\href{https://github.com/ChristianKudeba/alpoge-rh-verification/blob/alpoge-source-7cbf7a0c/src/Development/AutomorphicArithmetic/AutomorphicVarieties/PrimitiveFrobeniusFixed/DerInv.lean}{invariance modulo powers},
\href{https://github.com/ChristianKudeba/alpoge-rh-verification/blob/alpoge-source-7cbf7a0c/src/Development/AutomorphicArithmetic/AutomorphicVarieties/PrimitiveFrobeniusFixed/KolHom.lean}{Kummer classes},
\href{https://github.com/ChristianKudeba/alpoge-rh-verification/blob/alpoge-source-7cbf7a0c/src/Development/AutomorphicArithmetic/AutomorphicVarieties/PrimitiveFrobeniusFixed/KolInit.lean}{initial local computation},
\href{https://github.com/ChristianKudeba/alpoge-rh-verification/blob/alpoge-source-7cbf7a0c/src/Development/AutomorphicArithmetic/AutomorphicVarieties/PrimitiveFrobeniusFixed/KolStep.lean}{transition local computation},
\href{https://github.com/ChristianKudeba/alpoge-rh-verification/blob/alpoge-source-7cbf7a0c/src/Development/AutomorphicArithmetic/AutomorphicVarieties/PrimitiveFrobeniusFixed/NuSysEx.lean}{projection to the dual character and the two units}.}

\subsection{One prime for the two Frobenius conditions}
\label{sec:minus-frobenius}

Write \(U=G_{\mathbb Q(\mu_P)}\), \(W=\operatorname{Frac}(\mathcal O)/\mathcal O\),
and let \(\varpi\) be a uniformizer of \(\mathcal O\).
Let \(c,\nu:U\to W[M]\) be continuous additive homomorphisms such
that \(c(gug^{-1})=\alpha(g)c(u)\) and
\(\nu(gug^{-1})=\beta(g)\nu(u)\) for \(u\in U\) and
\(g\in G_{\mathbb Q}\). Assume \(c\) is nonzero on \(U\).
Given any finite set of excluded primes, we shall select an
admissible \(q\) outside the set and a Frobenius \(\phi\in U\) at
a place above \(q\) for which
\[
    c(\phi)\ne0,\qquad \mathcal O\,\nu(\phi)=
          \mathcal O\,\nu(U).
\]
Here \(\mathcal O\,\nu(U)\) means the submodule generated by all
the values. In the application \(c\) is a nonzero normalized strict
cocycle restricted to \(U\). Indeed, if the restriction were zero,
the cocycle, with \(\alpha|_U=1\) on \(W[M]\), would descend to
the abelian cyclotomic quotient. Commuting \(g\) with the
chosen lift of inversion \(j\) there gives
\((1-\alpha(j))c(g)=(1-\alpha(g))c(j)=0\). As
\(\alpha(j)\equiv-1\bmod M\) and \(\ell\) is odd, this would
force \(c=0\).

Choose \(\rho^M=r\) and let \(V=\{u\in U:u\rho=\rho\}\).
The homomorphism \(u\mapsto u\rho/\rho\) embeds
\(U/V\) in \(\mu_M\). For \(g\in G_{\mathbb Q}\) and \(u\in U\),
the ratio \(g^{-1}\rho/\rho\) lies in \(\mu_M\) and is fixed by
\(u\). Therefore
\[
       (gug^{-1})\rho/\rho=g(u\rho/\rho).
\]
Consequently any \(\gamma\)-equivariant additive map
\(f:U\to W'\) into an \(M\)-torsion \(\mathcal O\)-module that
vanishes on \(V\) satisfies
\((\gamma(g)-\kappa(g))f(u)=0\) for every \(g,u\). This follows
by computing \(f(gug^{-1})\) on \(U/V\), where conjugation acts
by the exponent \(\kappa(g)\bmod M\).

Put \(B=\ker c\). The unit \(\alpha(\delta)-\kappa(\delta)\) and
\(c|_U\ne0\) show that \(V\) is not contained in \(B\).
For \(C=\mathcal O\,\nu(U)\ne0\), put
\(A=\{u\in U:\nu(u)\in\varpi C\}\). The values generate a finite
cyclic module: \(W[M]\) is a quotient of the discrete valuation
ring \(\mathcal O\) by a nonzero principal ideal, and all its
submodules are cyclic. The image modulo \(\varpi C\) is nonzero.
At a lift \(j\) of inversion we have
\(\beta(j)-\kappa(j)=\beta(j)(1-\alpha(j))\), a unit. Applying
the preceding calculation to \(f=\nu\bmod\varpi C\) proves that
\(V\) is not contained in \(A\). Each \(\phi\in U\setminus A\)
generates \(C\) under evaluation by \(\nu\). For \(C=0\) take
\(A=B\); the generator condition then holds automatically.

For the third subgroup put \(S=G_{\mathbb Q(\mu_{P\ell})}\).
Choose \(u\) acting on \(\mu_{P\ell}\) by \(1+P\), a unit
modulo \(P\ell\), and \(g\) acting by inversion. Then \(u\in U\),
and \(v=gug^{-1}u\) lies in \(V\) by the displayed radical formula.
Its exponent on \(\mu_{P\ell}\) is \((1+P)^2\); this differs
from \(1\bmod P\ell\) since \(\ell\nmid 2+P\).
Thus \(V\) is not contained in \(S\).

Here is the elementary simultaneous-avoidance argument.
A group cannot be covered by two proper subgroups: if neither
contains the other, the product of an element exclusive to each lies
in neither. If \(V\) were covered by \(A,B,S\), choose
\(x\in V\setminus(B\cup S)\) and \(y\in V\setminus(A\cup S)\).
Then \(x\in A,\ y\in B,\ xy\in S\). The element \(x^2y\) cannot
lie in \(A\), since \(y\notin A\), or in \(S\), since \(xy\in S\)
and \(x\notin S\). If \(x^2y\in B\), then \(x^2\in B\).
Also \(x^M\in B\) since \(M c(x)=0\). As \(M\) is odd, these
two powers imply \(x\in B\), a contradiction. Hence
\(V\setminus(A\cup B\cup S)\) is nonempty. This set is open:
the stabilizers are open and the kernels defining \(A,B\) are
open and closed by continuity to finite discrete groups.

Choose any coset of an open normal subgroup of \(G_{\mathbb Q}\)
contained in this set. Chebotarev's theorem in the corresponding
finite Galois extension gives an unramified \(q\) outside the
excluded set and the divisors of \(P\ell r\) with the conjugacy
class of that coset as Frobenius. A choice of place above \(q\)
realizes the prescribed representative, and an arithmetic Frobenius
lift \(\phi\) then belongs to our open set. The cyclotomic action
gives \(q\equiv1\bmod P\) and \(q\not\equiv1\bmod P\ell\).
The element \(\rho\) is a unit at this place since \(q\ne r\);
as \(\phi\rho=\rho\), its reduction is fixed by \(q\)-power
Frobenius and hence lies in \(\mathbb F_q^\times\). It follows
that \(r^{(q-1)/M}\equiv1\bmod q\). Because
\(v_\ell(q-1)=m\), the order of \(r\bmod q\) is prime to
\(\ell\). These are precisely the admissibility conditions used
in the derivative relations. The exclusions from \(A\) and \(B\)
give the two simultaneous evaluations.
\footnote{\href{https://github.com/ChristianKudeba/alpoge-rh-verification/blob/alpoge-source-7cbf7a0c/src/Development/AutomorphicArithmetic/AutomorphicVarieties/PrimitiveFrobeniusFixed/AddPowNsmul.lean}{Selection of the auxiliary prime and simultaneous avoidance},
\href{https://github.com/ChristianKudeba/alpoge-rh-verification/blob/alpoge-source-7cbf7a0c/src/Development/AutomorphicArithmetic/AutomorphicVarieties/PrimitiveFrobeniusFixed/RadEqv.lean}{action on the radical quotient},
\href{https://github.com/ChristianKudeba/alpoge-rh-verification/blob/alpoge-source-7cbf7a0c/src/Development/AutomorphicArithmetic/AutomorphicVarieties/PrimitiveFrobeniusFixed/ChebQ.lean}{Chebotarev with a finite set of exclusions}.}
